% section: 漸化式が組合せ論と確率へ進む

\section{漸化式が組合せ論と確率へ進む}

\subsection{パスカルの三角形}

二項係数は、
\[
  \binom{n+1}{k}
  =
  \binom{n}{k-1}+\binom{n}{k}
\]
を満たす。
これは、$(n+1)$ 個から $k$ 個を選ぶとき、特別な一個を選ぶ場合と選ばない場合に分けたものである。

二項係数は、二次元の漸化式である。
一方向に並んだ数列だけでなく、格子上の位置 $(n,k)$ を動く規則も漸化式と考えられる。

この漸化式を繰り返すと、
\[
  \binom{n}{0},\binom{n}{1},\ldots,\binom{n}{n}
\]
ができ、パスカルの三角形になる。
さらに、
\[
  \sum_{k=0}^{n}\binom{n}{k}x^k=(1+x)^n
\]
という母関数的な式も現れる。

漸化式、組合せ、二項定理は、別々の章の内容ではない。
「一つの対象を、最後の選択によって場合分けする」という同じ考え方を持っている。

\subsection{確率遷移と行列}

二つの状態 $A,B$ を行き来する確率を考える。
時刻 $n$ に状態 $A,B$ にいる確率を $u_n,v_n$ とし、
\[
  \begin{pmatrix}u_{n+1}\\v_{n+1}\end{pmatrix}
  =
  \begin{pmatrix}p&q\\1-p&1-q\end{pmatrix}
  \begin{pmatrix}u_n\\v_n\end{pmatrix}
\]
とする。
各列または各行の和が $1$ になる行列を、確率の遷移行列として使う。

これは本編で見た連立漸化式と同じ形である。
行列の固有値のうち、$1$ 以外のものの絶対値が $1$ より小さければ、初期状態の影響が消え、定常的な確率分布へ近づく。

数列の極限、行列の固有値、確率過程の定常分布が、同じ線形代数の言葉で説明される。
漸化式は、確率を扱うための最も基本的な時間発展の記法でもある。

